\chapter{モールス信号}
% full version: chapters/04_code.tex

モールス信号には2つの記号、トンとツーがあり、その間に間（ま）がある。ニューロンのアルファベットはもっと貧しい。記号は1つ、短いカチッだけだ。つまりメッセージはすべて間に書き込まれている。では、ニューロンはどんな言葉で話しているのだろうか。

\section*{カチッでどれだけ語れるか}

通信エンジニアのように始めよう。こうした通信路には、そもそもどれだけの情報が入るのか。受け手がカチッの到着時刻を1ミリ秒の精度で区別できるなら、1回のカチッで約8ビットを運べる。受け手が1秒間に届いたカチッの数を数えるだけなら、同じ流れから取り出せる情報は数十分の1になる。通信路の容量のほとんどは\emph{時間}の中にある。

\pencilfig{4}{\pencilart{4}}{上はモールス信号、下はニューロン。記号はカチッ1つだけで、メッセージはすべて間に書き込まれている。}

ただし通信路にはノイズがある。意外なことに、ニューロンそのものは正確だ。同じ変動する電流を何度も与えると、ミリ秒の精度で同じ瞬間にカチッと鳴る。ノイズを出すのは接続部である。配線の端までたどり着いたカチッの半分以上は、ただ失われる。物質が毎回放出されるわけではないからだ。生きた大脳皮質では、カチッの流れはほとんどランダムに見える。これはノイズなのか。それとも、まだ読み方がわからないメッセージなのか。よく圧縮されたメッセージは、部外者にはランダムと区別がつかない。だからこの問いは未解決である。

\section*{遅いが確実}

最も単純な符号は、カチッを数えることだ。頻繁なほど数が大きい。この符号は損失にも揺らぎにも強い。だが遅い。頻度を10パーセントの精度で知るには100回のカチッが必要で、それには数秒かかる。ところがあなたは、一瞬だけ見せられた絵の中の動物をおよそ0.15秒で見分ける。その間に信号は約10段の処理を通る。各段で、各ニューロンはせいぜい1回しか鳴れない。

解決策は、たくさんのニューロンを同時に使うことだ。1人の聞き手が1分かけてカチッを数えるのではなく、100人の聞き手が一瞬で数える。ただし、ここにも落とし穴がある。隣り合うニューロンのノイズは、群衆の中のうわさのように少し似ている。数十個の隣で平均をとると効果はあるが、それ以上はほとんど効かなくなる。

\pencilfig{122}{%
% error of the group-average rate relative to one neuron, sqrt(c + (1-c)/N): c = 0.1 (neighbours' noise
% is slightly alike; full edition, chapter 4) and c = 0 (independent noise). x = 0.8 log2 N, y = 2.2 * error
\draw[pen] (0,0) -- (5.9,0); \draw[pen] (0,0) -- (0,2.45);
\foreach \x/\n in {0/1, 2.66/10, 5.31/100}{\draw[pcGray, line width=0.5pt] (\x,0) -- (\x,-0.08); \node[lbl, anchor=north] at (\x,-0.08) {\n};}
\node[lbl, anchor=north] at (2.95,-0.38) {平均するニューロンの数};
\node[lbl, anchor=south, rotate=90] at (-0.1,1.2) {誤差};
\draw[pcGray, densely dashed, line width=0.5pt] (0,0.70) -- (5.6,0.70);
\draw[penlite=pcBlue] plot[smooth] coordinates {(0.00,2.20) (0.47,1.80) (0.80,1.56) (1.27,1.27) (1.60,1.10) (2.07,0.90) (2.40,0.78) (2.87,0.64) (3.20,0.55) (3.67,0.45) (4.00,0.39) (4.47,0.32) (4.80,0.28) (5.27,0.22) (5.60,0.19)};
\draw[penlite=pcRed] plot[smooth] coordinates {(0.00,2.20) (0.47,1.84) (0.80,1.63) (1.27,1.39) (1.60,1.25) (2.07,1.10) (2.40,1.01) (2.87,0.92) (3.20,0.87) (3.67,0.82) (4.00,0.79) (4.47,0.76) (4.80,0.74) (5.27,0.73) (5.60,0.72)};
\node[lbl, text=pcRed, anchor=west, align=left] at (5.7,0.74) {似たノイズ\\（皮質）};
\node[lbl, text=pcBlue, anchor=west] at (5.7,0.19) {独立なノイズ};
\node[lbl, anchor=west] at (0.9,2.15) {ニューロン1個};
\node[lbl, anchor=north west] at (0.05,0.66) {3分の1};
}{ニューロン群で平均した頻度の誤差。ノイズが独立なら、誤差はさらに下がり続けるはずだ。隣どうしの似たノイズが誤差をおよそ3分の1で止め、その限界には数十個のニューロンですでに達する。}

\section*{速い：誰が一番か}

2つ目の解決策は、数を時間で書くことだ。強い信号はニューロンを早く鳴らし、弱い信号は遅く鳴らす。たった1回のカチッがすでに数を運ぶ。だが、時計がないなら何から時間を測ればいいのか。ニューロンどうしを比べればよい。誰が1番に鳴り、誰が2番に鳴ったか。10個のニューロンのカチッの到着順は、20ビット以上を運ぶ。さらに、共通の一斉発火が開始の目印になる。モールス信号で単語の間に置く長い間のようなものだ。

\pencilfig{121}{%
% latency code: one click per neuron; the stronger the input, the earlier the click
\foreach \y/\t/\l/\k in {1.5/1.1/{強い入力}/1, 0.95/2.4/{中くらい}/2, 0.4/4.2/{弱い}/3}{
  \draw[pen] (0.4,\y) -- (5.1,\y);
  \draw[pcBlue, line width=0.9pt] (\t,\y) -- (\t,\y+0.42);
  \node[lbl, anchor=east] at (0.3,\y+0.1) {\l};
  \draw[dotpen=pcAmber, wash=pcAmber] (5.6,\y+0.16) circle (0.17);
  \node[font=\scriptsize, text=pcAmber!60!black] at (5.6,\y+0.16) {\k};}
\draw[pcGray, densely dashed, line width=0.5pt] (0.55,0.25) -- (0.55,2.1);
\node[lbl, anchor=south] at (0.55,2.1) {刺激};
\node[lbl, text=pcAmber!70!black, anchor=south] at (5.6,2.1) {順番};
\draw[arr] (0.7,0.05) -- (4.9,0.05); \node[lbl, anchor=north] at (2.8,0.02) {時間};
}{強い入力は早いカチッ、弱い入力は遅いカチッ。開始の瞬間がわからなくても、カチッの順番がどの入力が強いかを教えてくれる。}

\section*{符号を選ぶのは聞き手}

同じカチッの列が、頻度も、時刻も、順番も運んでいる。そのどれが読み取られるかは受け手が決める。バケツの「穴が広い」ニューロンはゆっくり水をため、頻度だけを聞く。すぐ忘れるニューロンは同時発生だけを聞く。そして前に見たように、ブレーキはニューロンをある動作モードから別のモードへ切り替えられる。

接続部さえ信号をふるいにかける。ある接続部はすぐ疲れるので、頻度ではなくその\emph{変化}を伝える。逆に、数回続けてのカチッのバーストをよく通す接続部もある。ニューロンに「ツー」ができるわけだ。その間、抑制性ニューロンは句読点を打つ。流れを窓に区切り、うるさすぎるときは音量を絞る。

\section*{倹約な言葉}

最後に、この言葉は高くつく。カチッ1回ごとにエネルギーがかかり、脳全体の消費は約20ワットである。このワット数を全ニューロンで割ると、平均でニューロン1個あたり毎秒1回程度のカチッになる。ほとんどのニューロンはほとんどの時間黙っている。脳の符号はけちでなければならない。

モールス信号では、英語で最も頻繁な文字Eはトン1つだ。頻繁なものは短く符号化する。ニューロンも似ている。予想どおりのことにはカチッをほとんど使わない。一定の刺激はしだいに応答を引き起こさなくなり、センサーは変化を報告し、ドーパミンは驚きだけを報告する。脳はカチッをニュースに使うのである。

\fullversion{4}
